\documentclass[12pt, a4paper]{article}

% ── Paquetes ────────────────────────────────────────────────────────────────
\usepackage[utf8]{inputenc}
\usepackage[T1]{fontenc}
\usepackage[spanish, mexico]{babel}
\usepackage{amsmath, amssymb, amsthm}
\usepackage{geometry}
\usepackage{graphicx}
\usepackage{xcolor}
\usepackage{mdframed}
\usepackage{enumitem}
\usepackage{booktabs}
\usepackage{hyperref}

% ── Geometría ───────────────────────────────────────────────────────────────
\geometry{
  top    = 2.5cm,
  bottom = 2.5cm,
  left   = 3.0cm,
  right  = 3.0cm
}

% ── Colores ─────────────────────────────────────────────────────────────────
\definecolor{azulviz}{RGB}{37, 99, 235}
\definecolor{naranjviz}{RGB}{251, 146, 60}
\definecolor{verdviz}{RGB}{52, 211, 153}
\definecolor{amarillviz}{RGB}{251, 191, 36}
\definecolor{fondocaja}{RGB}{240, 245, 255}
\definecolor{bordecruz}{RGB}{37, 99, 235}
\definecolor{fondoadv}{RGB}{255, 248, 230}
\definecolor{bordeadv}{RGB}{251, 146, 60}
\definecolor{fondoteo}{RGB}{235, 250, 242}
\definecolor{bordeteo}{RGB}{52, 153, 110}

% ── Entornos ─────────────────────────────────────────────────────────────────
\mdfdefinestyle{vizverde}{
  backgroundcolor = fondoteo,
  linecolor       = bordeteo,
  linewidth       = 1.5pt,
  leftmargin      = 0pt,
  rightmargin     = 0pt,
  innertopmargin  = 8pt,
  innerbottommargin = 8pt,
  innerleftmargin = 12pt,
  innerrightmargin= 12pt,
  roundcorner     = 4pt
}

\mdfdefinestyle{viznaranja}{
  backgroundcolor = fondoadv,
  linecolor       = bordeadv,
  linewidth       = 1.5pt,
  leftmargin      = 0pt,
  rightmargin     = 0pt,
  innertopmargin  = 8pt,
  innerbottommargin = 8pt,
  innerleftmargin = 12pt,
  innerrightmargin= 12pt,
  roundcorner     = 4pt
}

\mdfdefinestyle{vizazul}{
  backgroundcolor = fondocaja,
  linecolor       = bordecruz,
  linewidth       = 1.5pt,
  leftmargin      = 0pt,
  rightmargin     = 0pt,
  innertopmargin  = 8pt,
  innerbottommargin = 8pt,
  innerleftmargin = 12pt,
  innerrightmargin= 12pt,
  roundcorner     = 4pt
}

\newenvironment{viznota}[1]{%
  \begin{mdframed}[style=vizverde]
  \textbf{\color{bordeteo}En GradienViz \textemdash\ #1.}\quad
}{%
  \end{mdframed}
}

\newenvironment{vizadv}{%
  \begin{mdframed}[style=viznaranja]
  \textbf{\color{bordeadv}Atención:}\quad
}{%
  \end{mdframed}
}

\newenvironment{vizdef}{%
  \begin{mdframed}[style=vizazul]
}{%
  \end{mdframed}
}

% ── Comandos matemáticos ─────────────────────────────────────────────────────
\newcommand{\R}{\mathbb{R}}
\newcommand{\norm}[1]{\left\lVert #1 \right\rVert}
\newcommand{\grad}{\nabla}
\newcommand{\mstar}{m^{*}}
\newcommand{\bstar}{b^{*}}
\newcommand{\Jfun}{J(m,b)}
\newcommand{\XtX}{\mathbf{X}^{\top}\mathbf{X}}
\newcommand{\T}{{\top}}

% ── Hyperref ─────────────────────────────────────────────────────────────────
\hypersetup{
  colorlinks = true,
  linkcolor  = azulviz,
  urlcolor   = azulviz,
  citecolor  = azulviz
}

% ── Metadatos ────────────────────────────────────────────────────────────────
\title{
  {\Large\bfseries Descenso por gradiente en regresión lineal}\\[0.4em]
  {\large Del visualizador a la red neuronal}\\[0.8em]
  {\normalsize Notas de lectura para \textit{GradienViz}}
}
\author{
  Departamento Académico de Sistemas Computacionales\\
  Universidad Autónoma de Baja California Sur
}
\date{2026}

% ─────────────────────────────────────────────────────────────────────────────
\begin{document}
\maketitle
\tableofcontents
\vspace{1em}

\begin{vizadv}
  Este documento asume familiaridad con cálculo diferencial multivariable,
  álgebra lineal elemental y la fórmula de regresión lineal simple.
  Se recomienda tener abierto \textbf{GradienViz} mientras se lee.
\end{vizadv}

% =============================================================================
\section{Punto de partida: lo que ya viste}
% =============================================================================

Cuando abres GradienViz y colocas algunos puntos en el Panel~1, la herramienta
calcula automáticamente dos números: la pendiente~$\mstar$ y el intercepto~$\bstar$
de la recta que mejor se ajusta a esos datos, y los marca con una cruz amarilla~$\oplus$
en el Panel~2.

Esos dos números son la solución al siguiente problema:

\begin{vizdef}
\textbf{Problema de regresión lineal simple.} Dados $n$ pares de observaciones
$(x_i, y_i)$ con $i = 1, \ldots, n$, encontrar los valores $m, b \in \R$ que
minimizan el error cuadrático medio:
\begin{equation}\label{eq:J}
  \Jfun = \frac{1}{n}\sum_{i=1}^{n}\bigl(m x_i + b - y_i\bigr)^2.
\end{equation}
\end{vizdef}

Ya conoces la solución analítica. Definiendo las sumas
\[
  S_x = \sum x_i, \quad
  S_y = \sum y_i, \quad
  S_{xy} = \sum x_i y_i, \quad
  S_{x^2} = \sum x_i^2,
\]
se obtiene directamente:
\begin{equation}\label{eq:exacta}
  \mstar = \frac{n S_{xy} - S_x S_y}{n S_{x^2} - S_x^2},
  \qquad
  \bstar = \frac{S_y - \mstar S_x}{n}.
\end{equation}

La pregunta natural es: \textbf{si ya existe esta fórmula, ¿para qué simular
un algoritmo?} La respuesta a esa pregunta es el núcleo de este documento
y la justificación completa del descenso por gradiente como método general
de optimización en aprendizaje automático.

\begin{viznota}{Panel~1}
La recta gris punteada que siempre aparece es $y = \mstar x + \bstar$,
calculada con \eqref{eq:exacta}. La recta naranja, durante la animación,
corresponde a una estimación provisional que el algoritmo va corrigiendo
paso a paso. Al converger, ambas rectas coinciden.
\end{viznota}

% =============================================================================
\section{La función de costo como superficie}
% =============================================================================

La expresión~\eqref{eq:J} define una función $J: \R^2 \to \R$.
Para un conjunto de datos fijo, cada par $(m, b)$ produce un valor de costo~$J$.
Si graficamos este valor sobre el plano $(m, b)$, obtenemos una superficie.

Expandiendo el cuadrado en~\eqref{eq:J}:
\begin{align}
  \Jfun &= \frac{1}{n}\sum_{i=1}^{n}
    \bigl(m^2 x_i^2 + 2mb x_i - 2m x_i y_i + b^2 - 2b y_i + y_i^2\bigr) \notag\\
  &= \frac{S_{x^2}}{n}\,m^2 + \frac{2S_x}{n}\,mb + b^2
     - \frac{2S_{xy}}{n}\,m - \frac{2S_y}{n}\,b + \frac{S_{y^2}}{n}.
  \label{eq:Jexpandida}
\end{align}

Esta expresión es un \textbf{polinomio de grado~2} en $m$ y $b$, con
todos los términos de grado par o cruzados. En términos matriciales,
si se define el vector de parámetros $\boldsymbol{\theta} = (m, b)^\T$,
la función se puede escribir como
\begin{equation}\label{eq:Jmatrix}
  J(\boldsymbol{\theta}) = \frac{1}{n}
  \bigl(\mathbf{X}\boldsymbol{\theta} - \mathbf{y}\bigr)^\T
  \bigl(\mathbf{X}\boldsymbol{\theta} - \mathbf{y}\bigr),
\end{equation}
donde $\mathbf{X} \in \R^{n\times 2}$ es la \emph{matriz de diseño}:
\[
  \mathbf{X} = \begin{pmatrix} x_1 & 1 \\ x_2 & 1 \\ \vdots & \vdots \\ x_n & 1 \end{pmatrix},
  \qquad
  \mathbf{y} = \begin{pmatrix} y_1 \\ y_2 \\ \vdots \\ y_n \end{pmatrix}.
\]

\subsection{Convexidad}

La función~\eqref{eq:Jmatrix} es \textbf{convexa}: su Hessiana
\begin{equation}\label{eq:hessiana}
  H = \nabla^2 J = \frac{2}{n}\XtX
\end{equation}
es semidefinida positiva para cualquier $\mathbf{X}$ (pues $\mathbf{v}^\T \XtX \mathbf{v}
= \norm{\mathbf{X}\mathbf{v}}^2 \geq 0$ para todo $\mathbf{v}$).
Cuando los puntos no son colineales verticalmente, $H$ es \emph{definida} positiva,
lo que garantiza que $J$ tiene un único mínimo global: la solución
$(\mstar, \bstar)$.

Esta propiedad tiene una consecuencia visual directa.

\begin{viznota}{Panel~2, mapa de calor}
El colormap viridis que se muestra es la función $\Jfun$: azul oscuro cerca
del mínimo, amarillo lejos de él. Que el mapa forme siempre un único
\emph{valle} sin mesetas ni mínimos secundarios es la expresión visual
de la convexidad de~$J$.
\end{viznota}

% =============================================================================
\section{El gradiente: las flechas del Panel~2}
% =============================================================================

Para minimizar $J$, necesitamos saber en qué dirección decrece más rápidamente
en cada punto del plano $(m, b)$. Esa información la proporciona el gradiente.

Derivando \eqref{eq:Jexpandida} con respecto a cada parámetro:
\begin{align}
  \frac{\partial J}{\partial m} &=
    \frac{2}{n}\sum_{i=1}^{n}(m x_i + b - y_i)\,x_i,
  \label{eq:gm}\\[4pt]
  \frac{\partial J}{\partial b} &=
    \frac{2}{n}\sum_{i=1}^{n}(m x_i + b - y_i).
  \label{eq:gb}
\end{align}

El vector gradiente en el punto $(m, b)$ es:
\[
  \grad J(m, b) =
  \begin{pmatrix} \partial J/\partial m \\ \partial J/\partial b \end{pmatrix} \in \R^2.
\]

El gradiente tiene dos propiedades geométricas fundamentales:
\begin{enumerate}[label=(\roman*)]
  \item Apunta en la dirección de \textbf{mayor crecimiento} de $J$.
  \item Es \textbf{perpendicular} a las curvas de nivel de $J$.
\end{enumerate}

Por lo tanto, $-\grad J(m,b)$ apunta en la dirección de mayor \emph{descenso}.

\begin{viznota}{Panel~2, campo vectorial}
Cada flecha pequeña del campo es el vector $-\grad J$ en ese punto de la cuadrícula,
normalizado para que su longitud sea proporcional a $\norm{\grad J}$.
Las flechas apuntan siempre hacia el mínimo, convergiendo en~$\oplus$.
Cerca del óptimo las flechas se acortan: el gradiente tiende a cero conforme
$\boldsymbol{\theta} \to (\mstar, \bstar)$.
\end{viznota}

% =============================================================================
\section{El algoritmo de descenso por gradiente}
% =============================================================================

La idea del descenso por gradiente es simple: desde un punto inicial arbitrario
$\boldsymbol{\theta}^{(0)} = (m_0, b_0)$, moverse en la dirección del gradiente
negativo con un paso de tamaño controlado.

\begin{vizdef}
\textbf{Iteración de descenso por gradiente.} Dado $\boldsymbol{\theta}^{(t)} = (m^{(t)}, b^{(t)})$,
calcular:
\begin{equation}\label{eq:update}
  \boldsymbol{\theta}^{(t+1)} = \boldsymbol{\theta}^{(t)} - \eta\,\grad J\!\left(\boldsymbol{\theta}^{(t)}\right),
\end{equation}
o explícitamente:
\begin{align*}
  m^{(t+1)} &= m^{(t)} - \eta\,\frac{\partial J}{\partial m}\biggr|_{(m^{(t)},\,b^{(t)})}, \\[4pt]
  b^{(t+1)} &= b^{(t)} - \eta\,\frac{\partial J}{\partial b}\biggr|_{(m^{(t)},\,b^{(t)})}.
\end{align*}
La constante $\eta > 0$ se llama \textbf{tasa de aprendizaje} (\emph{learning rate}).
\end{vizdef}

\subsection{El rol de la tasa de aprendizaje \texorpdfstring{$\eta$}{eta}}

La tasa de aprendizaje $\eta$ controla el tamaño de cada paso y determina
el comportamiento del algoritmo:

\begin{center}
\begin{tabular}{lll}
  \toprule
  Régimen & Comportamiento & Causa \\
  \midrule
  $\eta$ muy pequeño & Convergencia lenta pero segura & Pasos insuficientes \\
  $\eta$ moderado    & Convergencia suave hacia $(\mstar,\bstar)$ & Equilibrio \\
  $\eta$ grande      & Oscilaciones o divergencia & El paso supera el mínimo \\
  \bottomrule
\end{tabular}
\end{center}

Para una función cuadrática como~\eqref{eq:Jexpandida}, puede demostrarse que
la condición suficiente de convergencia es:
\begin{equation}\label{eq:condeta}
  \eta < \frac{2}{\lambda_{\max}(H)} = \frac{n}{\lambda_{\max}(\XtX)},
\end{equation}
donde $\lambda_{\max}$ denota el eigenvalor máximo de la Hessiana.
Cuando $\eta$ supera este umbral, la iteración diverge.

\begin{viznota}{slider $\eta$}
El rango $[0.001, 0.5]$ en escala logarítmica permite explorar
los tres regímenes. Con datos bien distribuidos en $x$, el umbral
de divergencia suele estar alrededor de $\eta \approx 0.3$\textendash$0.4$.
Con datos mal condicionados, ese umbral puede ser mucho menor.
\end{viznota}

\subsection{Criterio de convergencia}

Se dice que una trayectoria converge cuando el gradiente se vuelve
suficientemente pequeño:
\[
  \norm{\grad J\!\left(\boldsymbol{\theta}^{(t)}\right)} < \varepsilon,
\]
con $\varepsilon = 10^{-4}$ en GradienViz. Este criterio es más robusto que
verificar el cambio en $J$, pues detecta correctamente los puntos estacionarios
incluso en superficies muy planas.

% =============================================================================
\section{Una simulación ilustrativa}
% =============================================================================

En este punto es necesario ser explícito sobre el carácter del ejercicio:

\begin{vizadv}
  En regresión lineal, el descenso por gradiente es \textbf{innecesario
  desde el punto de vista computacional}: la solución exacta~\eqref{eq:exacta}
  se calcula en tiempo $O(n)$ con aritmética flotante estándar.
  GradienViz es una herramienta de \textbf{visualización pedagógica}, no
  un método de cálculo. Su propósito es hacer visible el comportamiento
  del algoritmo en un caso donde la respuesta correcta se conoce y puede
  usarse como referencia (la cruz~$\oplus$).
\end{vizadv}

Esto tiene una consecuencia importante: el visualizador permite verificar
directamente si el algoritmo llegó al lugar correcto, comparando la posición
final de cada trayectoria con $(\mstar, \bstar)$. En el caso real\textemdash
una red neuronal con miles de parámetros\textemdash esa verificación directa
no es posible.

% =============================================================================
\section{Condicionamiento numérico}
% =============================================================================

\subsection{La matriz \texorpdfstring{$\XtX$}{X'X} y su geometría}

La Hessiana~\eqref{eq:hessiana} es proporcional a $\XtX$, una matriz $2\times 2$
simétrica y semidefinida positiva. Sus eigenvalores $\lambda_1 \geq \lambda_2 \geq 0$
determinan la forma geométrica del bowl de~$J$:

\begin{itemize}
  \item Si $\lambda_1 \approx \lambda_2$: el bowl es casi circular, el gradiente
    apunta casi directamente al mínimo, y el descenso converge en pocas iteraciones.
  \item Si $\lambda_1 \gg \lambda_2$: el bowl es un elipsoide muy elongado.
    El gradiente apunta oblicuamente al mínimo, y el descenso toma un camino
    en zigzag que puede requerir miles de iteraciones.
\end{itemize}

\subsection{El número de condición}

La relación entre los eigenvalores extremos define el \textbf{número de condición}:
\begin{equation}\label{eq:kappa}
  \kappa(\XtX) = \frac{\lambda_1}{\lambda_2}.
\end{equation}

Para la matriz $2\times 2$
\[
  \XtX = \begin{pmatrix} S_{x^2} & S_x \\ S_x & n \end{pmatrix},
\]
los eigenvalores se calculan analíticamente:
\begin{align*}
  \text{tr} &= S_{x^2} + n, \\
  \det &= n\,S_{x^2} - S_x^2, \\
  \lambda_{1,2} &= \frac{\text{tr}}{2} \pm
    \sqrt{\left(\frac{\text{tr}}{2}\right)^2 - \det}.
\end{align*}

Un $\kappa$ grande indica que el problema está \emph{mal condicionado}:
pequeños cambios en los datos pueden producir grandes cambios en la solución,
y el algoritmo de descenso por gradiente converge lentamente.

\begin{viznota}{Panel~1, esquina superior derecha}
El valor $\kappa(\XtX)$ se actualiza cada vez que se modifica un punto.
Los colores indican el régimen: gris ($\kappa \leq 100$), amarillo ($\kappa \leq 1000$),
naranja ($\kappa > 1000$). Los casos con puntos agrupados en torno a un mismo
valor de~$x$ producen $\kappa$ extremadamente alto y trayectorias que no convergen
dentro del límite de 2000 iteraciones.
\end{viznota}

\subsection{¿Por qué los puntos verticales degradan el condicionamiento?}

Si todas las $x_i$ son iguales (digamos $x_i = c$), entonces:
\[
  \XtX = \begin{pmatrix} nc^2 & nc \\ nc & n \end{pmatrix}
       = n \begin{pmatrix} c^2 & c \\ c & 1 \end{pmatrix},
\]
cuyo determinante es $n(c^2 - c^2) = 0$. La matriz es singular: $\lambda_2 = 0$
y $\kappa = \infty$. Geométricamente, los datos no permiten determinar la pendiente
$m$ de ninguna recta\textemdash cualquier recta vertical pasa por todos los puntos.
Conforme los datos se aproximan a esta configuración, $\kappa$ crece sin límite.

% =============================================================================
\section{De la recta a la red neuronal}
% =============================================================================

Hemos establecido que, para regresión lineal simple, el descenso por gradiente
es un algoritmo funcionalmente redundante: la fórmula~\eqref{eq:exacta}
lo supera en todo sentido. ¿Por qué, entonces, el descenso por gradiente
es el algoritmo central del aprendizaje profundo?

\subsection{Escala: cuando la fórmula exacta deja de ser práctica}

La solución analítica para regresión lineal múltiple con $p$ predictores es:
\begin{equation}\label{eq:lsmulti}
  \boldsymbol{\theta}^* = (\mathbf{X}^\T \mathbf{X})^{-1}\mathbf{X}^\T \mathbf{y}.
\end{equation}

Calcular esta expresión requiere invertir una matriz $p \times p$.
Con $p = 2$ (como en GradienViz) esto es trivial.
Pero en un modelo de lenguaje grande, $p$ puede superar los $10^{11}$ parámetros.
Almacenar y factorizar una matriz de $10^{11} \times 10^{11}$ es computacionalmente
imposible con cualquier hardware concebible.

\subsection{No-convexidad: cuando la fórmula exacta no existe}

La razón más profunda es que los modelos no lineales, en particular las redes
neuronales, no tienen una función de costo cuadrática.

Considera una red neuronal simple con arquitectura $2 \to 4 \to 1$:
dos entradas, una capa oculta de cuatro neuronas con función de activación
no lineal $\sigma$ (por ejemplo, ReLU o sigmoide), y una salida escalar.
El número total de parámetros es:
\[
  (2 \times 4) + 4 + (4 \times 1) + 1 = 8 + 4 + 4 + 1 = 17.
\]

La salida de esta red para una entrada $\mathbf{x} \in \R^2$ es:
\[
  \hat{y}(\mathbf{x}; \boldsymbol{\theta}) =
    \mathbf{w}_2^\T\,\sigma\!\bigl(\mathbf{W}_1 \mathbf{x} + \mathbf{b}_1\bigr) + b_2,
\]
donde $\mathbf{W}_1 \in \R^{4\times 2}$, $\mathbf{b}_1 \in \R^4$, $\mathbf{w}_2 \in \R^4$
y $b_2 \in \R$ son los parámetros, y $\boldsymbol{\theta} \in \R^{17}$.

La función de costo correspondiente,
\[
  J(\boldsymbol{\theta}) = \frac{1}{n}\sum_{i=1}^{n}
    \bigl(\hat{y}(\mathbf{x}_i;\boldsymbol{\theta}) - y_i\bigr)^2,
\]
ya no es cuadrática en $\boldsymbol{\theta}$ debido a la no linealidad de $\sigma$.
Como consecuencia:

\begin{enumerate}[label=(\roman*)]
  \item La Hessiana $\nabla^2 J$ no es constante: depende de $\boldsymbol{\theta}$.
  \item $J$ no es convexa en general: puede tener múltiples mínimos locales,
    máximos locales y puntos de silla.
  \item No existe una fórmula cerrada equivalente a~\eqref{eq:lsmulti}.
\end{enumerate}

El descenso por gradiente\textemdash y sus variantes (Adam, RMSprop, Adagrad)---
es el único método de optimización práctico disponible.

\subsection{La geometría del espacio de parámetros: variedades singulares}

La geometría de la superficie $J(\boldsymbol{\theta})$ para una red neuronal
es cualitativamente diferente al bowl elipsoidal de GradienViz.

Para una red neuronal con función de activación suave, $J$ define una
\emph{variedad estadística} en el espacio de parámetros. A diferencia del
elipsoide, esta variedad tiene \textbf{singularidades}: regiones donde la
Hessiana es singular (no invertible), incluyendo valles planos extendidos,
crestas, y mesetas.

Estas singularidades surgen de manera inevitable por las simetrías del modelo.
Por ejemplo, en la red $2\to 4\to 1$, permutar las cuatro neuronas ocultas
produce exactamente la misma función de salida: hay $4! = 24$ configuraciones
de parámetros distintas que representan el mismo modelo. Cada punto del mínimo
global de $J$ tiene al menos 23 copias exactas en el espacio de parámetros,
todas igualmente válidas.

En la teoría de aprendizaje estadístico singular\footnote{%
  Watanabe, S. (2009). \textit{Algebraic Geometry and Statistical Learning Theory}.
  Cambridge University Press.}, se demuestra que la estructura de estas
singularidades determina propiedades del proceso de aprendizaje como
la generalización, el comportamiento en muestras pequeñas y la
complejidad efectiva del modelo\textemdash aspectos que no tienen
análogo en la geometría convexa de la regresión lineal.

\subsection{La regla de actualización es la misma}

A pesar de estas diferencias, la iteración fundamental es idéntica a~\eqref{eq:update}:
\[
  \boldsymbol{\theta}^{(t+1)} = \boldsymbol{\theta}^{(t)}
    - \eta\,\grad J\!\left(\boldsymbol{\theta}^{(t)}\right).
\]

Lo que cambia es cómo se calcula $\grad J$.
En regresión lineal, las derivadas parciales~\eqref{eq:gm}--\eqref{eq:gb}
se obtienen directamente. En una red neuronal, el cálculo del gradiente
se realiza mediante el algoritmo de \textbf{retropropagación}
(\emph{backpropagation}), que aplica la regla de la cadena de manera
sistemática a lo largo de la arquitectura de la red.

El visualizador GradienViz muestra, en su esencia más simple, el mismo
proceso que ocurre durante el entrenamiento de cualquier red neuronal:
una trayectoria en el espacio de parámetros que sigue la dirección del
gradiente negativo. En la red, ese espacio tiene 17, 17\,000 o $10^{11}$
dimensiones y no puede representarse visualmente; pero la idea subyacente
es la misma que se ve en el Panel~2.

% =============================================================================
\section*{Resumen}
% =============================================================================

\begin{center}
\begin{tabular}{lll}
  \toprule
  & \textbf{Regresión lineal (GradienViz)} & \textbf{Red neuronal $2\to 4\to 1$} \\
  \midrule
  Parámetros          & 2 $(m, b)$          & 17 \\
  Función de costo    & Cuadrática, convexa & No lineal, no convexa \\
  Mínimo global       & Único, analítico    & Múltiple, sin fórmula cerrada \\
  Hessiana            & Constante           & Depende de $\boldsymbol{\theta}$ \\
  Paisaje de $J$      & Elipsoide regular   & Variedad singular \\
  Solución exacta     & Sí \eqref{eq:exacta} & No \\
  Necesidad del alg.  & Ninguna (pedagógica) & Indispensable \\
  Regla de actualiz.  & \eqref{eq:update}   & \eqref{eq:update} (misma) \\
  Cálculo del gradiente & Fórmulas directas & Retropropagación \\
  \bottomrule
\end{tabular}
\end{center}

\bigskip

El descenso por gradiente es un algoritmo de propósito general cuya utilidad
real emerge cuando no existe solución analítica, cuando el espacio de parámetros
es de dimensión muy alta, o cuando la función de costo tiene una geometría
compleja. GradienViz ofrece la única situación en que ese algoritmo puede
verse funcionando en un espacio de dos dimensiones, con la solución correcta
siempre visible como referencia.

\end{document}
